% Created 2025-11-21 Fri 18:27
% Intended LaTeX compiler: pdflatex
\documentclass[10pt]{article}
\usepackage[utf8]{inputenc}
\usepackage[T1]{fontenc}
\usepackage{graphicx}
\usepackage{longtable}
\usepackage{wrapfig}
\usepackage{rotating}
\usepackage[normalem]{ulem}
\usepackage{amsmath}
\usepackage{amssymb}
\usepackage{capt-of}
\usepackage{hyperref}
\usepackage[top=0.1cm, bottom=1cm, left=1cm, right=1cm]{geometry}
\usepackage{sectsty}
\allsectionsfont{\centering\bfseries}
\pagenumbering{gobble}
\setlength{\parindent}{0pt}
\author{Manuel Gortarez}
\date{21 de Noviembre de 2025}
\title{Gototion}
\hypersetup{
 pdfauthor={Manuel Gortarez},
 pdftitle={Gototion},
 pdfkeywords={},
 pdfsubject={},
 pdfcreator={Emacs 30.2 (Org mode 9.7.11)}, 
 pdflang={English}}
\begin{document}

\maketitle
\section*{Arquitectura de Gototion}
\label{sec:org3d6ddcd}

\begin{itemize}
\item Telegram bot (input del usuario y recibir la confirmacion).
\item AWS lambda (serverless y orquesta el trabajo entre telegram, notion y el agente)
\item Amazon Bedrock (Para el LLM y el uso del agente)
\item Notion API
\end{itemize}

Como el agente solo se encargará de crear instancias en una base de datos ya establecidas
solo necesitamos la API de notion y no MCP. Para esto aprovecharemos los comandos del bot
de telegram para la creacion de las tareas.
\section*{Base de datos en Notion}
\label{sec:orgd90748d}

Los campos que contendra la base de datos en notion necesarios para la creacion de una tarea
son los siguientes:
\begin{itemize}
\item Nombre de la tarea (Investigar sobre funciones de activación)
\item Materia (Redes Neuronales)
\item Fecha de entrega (3 de Agosto de 2025)
\item Estado (No inicado, iniciado, finalizado)
\item Nivel de prioridad (baja, media, alta)
\item Nivel de esfuerzo (bajo, medio, alto)
\end{itemize}

En notion podemos poner funciones en la base de datos asi como automatizaciones que nos ayudaran
a quitarle trabajo al agente. Por ejemplo un comando que es bastante util es el de "dias faltantes"
que se puede calcular muy fácil. Además, podemos acomodar las tareas por medio de uno o mas campos.
La manera en que se hace es con una matriz:

\begin{center}
\begin{tabular}{rll}
Orden & Prioridad & Esfuerzo\\
\hline
1 & Alta & Bajo\\
2 & Alta & Medio\\
3 & Alta & Alto\\
4 & Media & Bajo\\
5 & Media & Medio\\
6 & Media & Alto\\
7 & Baja & Bajo\\
8 & Baja & Medio\\
9 & Baja & Alto\\
\end{tabular}
\end{center}

Ahora con los días faltantes de entrega podemos hacer lo siguiente, crear un campo que el usuario no
puede interactuar con el, lo llamaremos "urgencia", podemos hacer que tenga tres niveles:
\begin{itemize}
\item Urgente.
\item Proximo.
\item Sin prisa.
\end{itemize}

El algoritmo para ordenar las tareas se vuelve sencillo en base a la matriz y a la urgencia:
\begin{enumerate}
\item Ordenar tareas por nivel de urgencias.
\item Por cada nivel de urgencias ordenar en base a la matriz Prioridad x Esfuerzo.
\end{enumerate}

Con esto ya tenemos todo lo necesario en notion, al menos en la parte de la base de datos.
\section*{Bot de Telegram}
\label{sec:org87bb0d0}

Como primera aproximación no usarmemos todavia un LLM, lo que haremos es en base a texto
formateado y comandos de telegram es que actualizaremos la base de datos en notion. Algunos comandos
son:
\begin{itemize}
\item Credenciales: aquí pondremos las credenciales de los usuarios para poder interactuar son su base de
datos en notion.
\item Crear Tarea: para crear una tarea de la escuela.
\item Modificar Tarea: por si hay algo que modificar en una tarea. Aqui hay que ver como manejamos que el
usuario escoja que tarea modificar.
\item Eliminar Tarea
\item Lista de Tareas: el bot le proporciona al usuario las tareas que tiene pendientes ya ordenadas, de tal
manera que la primera tarea es la primera que el usuario debería de hacer en base a las métricas.
\end{itemize}

Por el moment como no usaremos todavia ni AWS ni un LLM una idea de creacion de tarea sería asi:
\begin{itemize}
\item \texttt{/CrearTarea} : funciones de activacion, redes neuronales, 03-08-25, Media, Baja.
\end{itemize}
\end{document}
